\section{Discusi\'on}
Los resultados obtenidos permiten caracterizar la capacidad de la plataforma ThingsBoard desplegada sobre la red IEEE~802.11ah (HaLow) bajo distintos patrones de carga. En primer lugar, la l\'inea base con \num{20} dispositivos confirma que, en un escenario de baja concurrencia, la infraestructura se comporta como se espera: las latencias se mantienen en el orden de pocos milisegundos, no se presentan p\'erdidas de mensajes y el enlace HaLow opera con un margen amplio respecto a su l\'imite de sensibilidad.

Al escalar progresivamente el n\'umero de dispositivos hasta \num{400} clientes (Prueba~2), el sistema mantiene una relaci\'on casi lineal entre el n\'umero de nodos y el \emph{throughput} agregado, alcanzando del orden de \SI{69}{\text{msg/s}} sin evidencias de saturaci\'on en el \emph{broker} o en el servidor de aplicaci\'on. La ausencia de p\'erdidas de mensajes y la estabilidad en el n\'umero de dispositivos conectados sugieren que, para la configuraci\'on de hardware utilizada, la plataforma puede operar con cientos de nodos simult\'aneos sin comprometer la entrega de telemetr\'ia.

Las pruebas de variaci\'on del intervalo de env\'io (Prueba~3) profundizan en este an\'alisis. Con \num{200} dispositivos, la reducci\'on del intervalo de \SI{10}{\second} a \SI{1}{\second} incrementa el \emph{throughput} en un orden de magnitud, pero las latencias promedio se mantienen en un rango razonable y las colas internas del broker no llegan a provocar p\'erdida de mensajes. El aumento del percentil 95 de latencia en los intervalos intermedios (especialmente \SI{2}{\second}) indica que existen reg\'imenes donde la combinaci\'on de concurrencia y frecuencia de env\'io exige m\'as al \emph{back-end}, aun sin llegar a una degradaci\'on cr\'itica.

En la Prueba~4 se observa que el tama\~no del \emph{payload} no afecta exclusivamente al ancho de banda de la red, sino tambi\'en al procesamiento aplicaci\'on--nivel. El perfil \textit{medium} exhibe las latencias m\'as altas, mientras que \textit{small} y \textit{large} logran latencias menores, lo que sugiere que, en esta plataforma, ciertos tamaños intermedios de mensaje pueden interactuar de forma menos eficiente con la serializaci\'on, el almacenamiento o la validaci\'on de telemetr\'ia.

La comparaci\'on entre env\'io directo y configuraciones con l\'imite de tasa (Prueba~5) evidencia un compromiso cl\'asico: el \emph{rate limit} protege al servidor reduciendo la cantidad total de mensajes y el \emph{throughput} efectivo, a costa de incrementar ligeramente la latencia promedio y de disminuir la resoluci\'on temporal de las mediciones. Desde la perspectiva de dise\~no de la soluci\'on, esto implica que la activaci\'on de un \emph{rate limiter} debe acompa\~narse de pol\'iticas claras sobre la frecuencia m\'inima exigida por la aplicaci\'on final.

Por \'ultimo, el escenario as\'incrono de la Prueba~6 confirma que la plataforma mantiene su estabilidad cuando distintos lotes de dispositivos se incorporan en momentos diferentes. Incluso en la configuraci\'on m\'as exigente, con \num{400} nodos y \emph{throughput} cercano a \SI{133}{\text{msg/s}}, no se observan p\'erdidas de mensajes ni desconexiones masivas, y la red HaLow conserva valores de RSSI y tasas de transmisi\'on coherentes con una operaci\'on segura. Esto sugiere que el l\'imite pr\'actico de la soluci\'on, en la configuraci\'on actual, se encuentra por encima del rango de carga estudiado, y que las primeras restricciones de escalabilidad emerger\'an probablemente en el \emph{back-end} de ThingsBoard o en el dimensionamiento del broker MQTT, m\'as que en la capa de acceso inal\'ambrico.
